\documentclass[a4paper,12pt]{article}
\usepackage[T2A]{fontenc}
\usepackage[utf8]{inputenc}
\usepackage[russian]{babel}
\usepackage{amsmath,amssymb,mathtools}
\renewcommand{\familydefault}{\sfdefault}
\usepackage{icomma}
\usepackage{geometry}
\geometry{left=18mm,right=18mm,top=18mm,bottom=19mm}
\usepackage{microtype}
\usepackage{tikz}
\usetikzlibrary{calc,arrows.meta,positioning,shapes.geometric}
\usepackage{tabularx,booktabs,longtable}
\usepackage{enumitem}
\usepackage{hyperref}
\usepackage{parskip}
\usepackage{fancyhdr}
\usepackage{setspace}
\usepackage{xcolor}
\usepackage{array}

\definecolor{accent}{HTML}{2B6F73}
\definecolor{darkaccent}{HTML}{19484B}
\definecolor{textgray}{HTML}{2E3335}
\definecolor{softgray}{HTML}{6D777A}
\definecolor{linegray}{HTML}{D8DEDF}

\hypersetup{colorlinks=true,linkcolor=darkaccent,urlcolor=darkaccent}
\onehalfspacing
\setlength{\parindent}{0pt}
\setlist[itemize]{leftmargin=6mm,itemsep=2pt,topsep=3pt}
\setlist[enumerate]{leftmargin=8mm,itemsep=4pt,topsep=4pt}
\setlength{\headheight}{14.5pt}
\pagestyle{fancy}
\fancyhf{}
\fancyhead[L]{\small\color{softgray}София Хоменко}
\fancyhead[R]{\small\color{softgray}29.09.26}
\fancyfoot[C]{\small\color{softgray}\thepage}
\renewcommand{\headrulewidth}{0pt}
\renewcommand{\footrulewidth}{0pt}

\newcommand{\topic}[1]{%
  \vspace{0.5em}\noindent{\Large\bfseries\color{darkaccent}#1}\par
  \vspace{0.25em}\noindent\color{linegray}\rule{\linewidth}{0.6pt}\color{textgray}\vspace{0.5em}
}
\newcommand{\key}[1]{\textbf{\color{darkaccent}#1}}
\newcommand{\formula}[1]{\[\boxed{\displaystyle #1}\]}
\newcommand{\warn}{\textbf{\color{darkaccent}Важно: }}

\begin{document}
\color{textgray}

% ---------- TITLE ----------
\thispagestyle{empty}
\begin{center}
\vspace*{25mm}
{\Large\color{softgray}Математика • ОГЭ}\par
\vspace{10mm}
{\Huge\bfseries\color{darkaccent}Чек-лист}\par
\vspace{5mm}
{\LARGE\bfseries Автомобильные шины, проценты\par и метод интервалов}\par
\vspace{15mm}
{\large София Хоменко}\par
\vspace{4mm}
{\normalsize 29 сентября 2026 года}\par
\vfill
{\small Лёвин Артём Александрович эксклюзивно для Софии Хоменко}\par
\vspace{12mm}
\begin{tikzpicture}[x=1cm,y=1cm]
  \draw[accent,line width=1.2pt]
  (0,0) .. controls (3.2,0.8) and (7.3,-0.8) .. (10.5,0)
        .. controls (13.1,0.6) and (15.6,0.2) .. (17.4,-0.15);
\end{tikzpicture}
\end{center}
\clearpage

% ---------- MAIN CHECKLIST ----------
\topic{1. Маркировка шины: что означает каждое число}

Маркировка вида \(215/50\,R17\) содержит три основных параметра.

\begin{itemize}
  \item \key{215} --- ширина шины \(B\) в миллиметрах.
  \item \key{50} --- высота боковины \(H\) в процентах от ширины \(B\).
  \item \key{R17} --- диаметр диска \(d\) в дюймах. Один дюйм равен \(25{,}4\) мм.
\end{itemize}

Отсюда:
\formula{H=B\cdot\frac{p}{100}},
\formula{d=R\cdot25{,}4}.

\textbf{Пример.} Для шины \(215/50\,R17\):
\[
B=215\text{ мм},\qquad H=215\cdot0{,}50=107{,}5\text{ мм},
\]
\[
d=17\cdot25{,}4=431{,}8\text{ мм}.
\]

\warn второе число в маркировке --- это \emph{процент}, а не высота в миллиметрах.

\topic{2. Диаметр всего колеса}

Полный диаметр колеса складывается из диаметра диска и двух боковин шины:
\formula{D=d+2H}.

Для заводской шины \(215/50\,R17\):
\[
D=431{,}8+2\cdot107{,}5=646{,}8\text{ мм}.
\]

Удобно сразу подставлять данные маркировки в одну формулу:
\formula{D=2B\cdot\frac{p}{100}+25{,}4R}.

\warn не путайте \(d\) --- диаметр \emph{диска} и \(D\) --- диаметр \emph{всего колеса}.

\topic{3. Как читать таблицу разрешённых размеров}

В первых заданиях ОГЭ часто нужно только аккуратно сопоставить строку и столбец таблицы.

\begin{enumerate}
  \item Найдите в условии заданный диаметр диска \(R\).
  \item Перейдите к соответствующему столбцу или строке таблицы.
  \item Смотрите только на разрешённые сочетания.
  \item Если спрашивают наибольшую ширину, сравнивайте значения \(B\) \emph{только внутри выбранного диаметра диска}.
\end{enumerate}

\warn наличие большего числа в соседнем столбце ещё не означает, что такой размер разрешён для нужного диска.

\topic{4. «На сколько увеличится?» и «на сколько процентов увеличится?»}

Если величина была \(A_{\text{стар}}\), а стала \(A_{\text{нов}}\), то абсолютное увеличение равно
\formula{\Delta A=A_{\text{нов}}-A_{\text{стар}}}.

Процентное увеличение:
\formula{q=\frac{A_{\text{нов}}-A_{\text{стар}}}{A_{\text{стар}}}\cdot100\%}.

\textbf{Главный ориентир:} в знаменателе стоит то, \key{что было изначально}.

\topic{5. Пробег за один оборот колеса}

За один полный оборот колесо проходит длину своей окружности:
\formula{L=\pi D}.

Если сравнивается процентное изменение пробега двух колёс, число \(\pi\) сокращается:
\[
\frac{L_{\text{нов}}-L_{\text{стар}}}{L_{\text{стар}}}\cdot100\%
=\frac{\pi D_{\text{нов}}-\pi D_{\text{стар}}}{\pi D_{\text{стар}}}\cdot100\%
=\frac{D_{\text{нов}}-D_{\text{стар}}}{D_{\text{стар}}}\cdot100\%.
\]

То есть для пятого пункта такого блока достаточно корректно найти два диаметра.

\warn округляйте только конечный результат, если условие отдельно не требует иного.

\clearpage

\topic{6. Метод интервалов: короткое повторение}

Для квадратного неравенства вида
\[
f(x)\ge 0,\qquad f(x)=ax^2+bx+c,
\]
удобен следующий маршрут.

\begin{enumerate}
  \item Перенесите всё в одну сторону: справа должен быть ноль.
  \item Решите уравнение \(f(x)=0\) и найдите корни.
  \item Отметьте корни на числовой прямой.
  \item Определите знак \(f(x)\) на одном из интервалов подстановкой удобной точки.
  \item При переходе через простой корень знак меняется.
  \item Выберите интервалы, соответствующие знаку неравенства.
  \item Для \(\le\) и \(\ge\) корни включаются; для \(<\) и \(>\) --- не включаются.
\end{enumerate}

\textbf{Пример.}
\[
x^2-3x-4\ge0.
\]
Корни: \(x=-1\) и \(x=4\). При \(x=0\) получаем \(-4<0\), поэтому в среднем интервале знак «минус», по краям --- «плюс».
\[
\boxed{x\in(-\infty;-1]\cup[4;+\infty)}.
\]

\warn сначала нормализуйте неравенство. Подстановка точки имеет смысл в выражение, у которого справа уже стоит ноль.

% ---------- FLOWCHART ----------
\clearpage
\thispagestyle{plain}
\begin{center}
{\Large\bfseries\color{darkaccent}Алгоритм: задача ОГЭ про автомобильные шины}\par
\vspace{8mm}
\end{center}

\tikzset{
  startstop/.style={ellipse,draw=accent,line width=0.9pt,align=center,minimum width=38mm,minimum height=10mm},
  process/.style={rounded corners=3pt,draw=accent,line width=0.9pt,align=center,text width=105mm,minimum height=12mm,inner sep=4mm},
  decision/.style={diamond,aspect=2.4,draw=accent,line width=0.9pt,align=center,text width=38mm,inner sep=1.8mm},
  arrow/.style={-{Latex[length=3mm]},line width=0.9pt,color=darkaccent}
}
\begin{center}
\begin{tikzpicture}[node distance=10mm]
  \node[startstop] (s) {Прочитать условие};
  \node[process,below=of s] (decode) {Расшифровать маркировку: \(B\), процент \(p\), диаметр диска \(R\)};
  \node[process,below=of decode] (calc) {Посчитать \(H=B\cdot p/100\), \(d=25{,}4R\), затем \(D=d+2H\)};
  \node[decision,below=12mm of calc] (q) {Что спрашивают?};
  \node[rounded corners=3pt,draw=accent,line width=0.9pt,align=center,text width=62mm,minimum height=14mm,below=16mm of q,xshift=-42mm] (abs) {Разницу в мм:\\[1mm]\(D_{\text{нов}}-D_{\text{стар}}\)};
  \node[rounded corners=3pt,draw=accent,line width=0.9pt,align=center,text width=62mm,minimum height=14mm,below=16mm of q,xshift=42mm] (pct) {Проценты / пробег:\\[1mm]\(\dfrac{D_{\text{нов}}-D_{\text{стар}}}{D_{\text{стар}}}\cdot100\%\)};
  \node[process,below=30mm of q] (round) {Округлить только в конце, если это требует условие};
  \node[startstop,below=of round] (f) {Записать ответ с нужной точностью};

  \draw[arrow] (s) -- (decode);
  \draw[arrow] (decode) -- (calc);
  \draw[arrow] (calc) -- (q);
  \draw[arrow] (q) -- node[above left,font=\small]{мм} (abs);
  \draw[arrow] (q) -- node[above right,font=\small]{\%} (pct);
  \draw[arrow] (abs.south) -- (round.north west);
  \draw[arrow] (pct.south) -- (round.north east);
  \draw[arrow] (round) -- (f);
\end{tikzpicture}
\end{center}
\clearpage

% ---------- TRAINING ----------
\topic{Тренировочный блок — Путь А}

\textbf{10 заданий.} Выполняйте вычисления вручную. В задачах на проценты округляйте ответ до десятых, если не сказано иное.

\begin{enumerate}
  \item Для шины \(205/55\,R16\) найдите ширину \(B\), высоту боковины \(H\), диаметр диска \(d\) в миллиметрах и полный диаметр колеса \(D\).

  \item Сколько миллиметров составляет высота боковины шины \(225/45\,R17\)?

  \item Найдите полный диаметр колеса с шиной \(195/65\,R15\).

  \item Заводское колесо имеет шину \(215/50\,R17\). Его заменили колесом с шиной \(205/60\,R16\). На сколько миллиметров изменился диаметр колеса?

  \item На сколько процентов увеличится пробег автомобиля за один оборот колеса при замене \(195/65\,R15\) на \(205/60\,R16\)?

  \item В таблице разрешённых размеров для диска \(R16\) указаны ширины \(195\), \(205\) и \(215\) мм; ширина \(225\) мм разрешена только для \(R17\) и \(R18\). Какова наибольшая разрешённая ширина шины для \(R16\)?

  \item У шины ширина \(205\) мм, диаметр диска \(R16\), а полный диаметр колеса равен \(631{,}9\) мм. Найдите второе число в маркировке шины.

  \item Решите неравенство
  \[
  x^2-5x-6<0.
  \]

  \item Решите неравенство
  \[
  2x^2-x-3\ge0.
  \]

  \item Решите неравенство
  \[
  x^2+x-12\le0.
  \]
\end{enumerate}

\clearpage

% ---------- ANSWERS ----------
\topic{Ответы}

\renewcommand{\arraystretch}{1.35}
\begin{longtable}{>{\raggedright\arraybackslash}p{11mm} >{\raggedright\arraybackslash}p{0.78\linewidth}}
\toprule
\textbf{№} & \textbf{Ответ} \\
\midrule
\endhead
1 & \(B=205\) мм; \(H=112{,}75\) мм; \(d=406{,}4\) мм; \(D=631{,}9\) мм. \\
2 & \(101{,}25\) мм. \\
3 & \(634{,}5\) мм. \\
4 & Увеличился на \(5{,}6\) мм. \\
5 & Примерно на \(2{,}8\%\). \\
6 & \(215\) мм. \\
7 & \(55\). Маркировка: \(205/55\,R16\). \\
8 & \(x\in(-1;6)\). \\
9 & \(x\in(-\infty;-1]\cup[\tfrac32;+\infty)\). \\
10 & \(x\in[-4;3]\). \\
\bottomrule
\end{longtable}

\vfill
\begin{center}
{\small\color{softgray}Лёвин Артём Александрович эксклюзивно для Софии Хоменко}
\end{center}

\end{document}
